\section{Il potenziale d'azione}

Nel neurone si distinguono 4 elementi con diversi ruoli funzionali:
\begin{itemize}
  \item \textbf{dendriti}: porzione di ingresso;
  \item \textbf{cono di emergenza}: elemento decisionale o \textit{trigger};
  \item \textbf{assone}: conduzione del messaggio all'interno della cellula;
  \item \textbf{sinapsi}: elemento di uscita e di comunicazione con altre cellule.
\end{itemize}

% ---------------------------------------------------------------------------
\subsection{Potenziale di membrana}

\begin{itemize}
  \item mantenuto costante dalla \textbf{pompa sodio/potassio} (estrude 3 $Na^+$ e introduce 2 $K^+$, con consumo di ATP);
  \item la costanza dei gradienti ionici assicura che il potenziale di membrana mantenga un valore costante ($\sim$$-70$\,mV) per l'intera durata della vita della cellula;
  \item in condizioni di riposo la cellula non si trova in uno stato di equilibrio, ma in uno \textbf{stato stazionario}: uno stato di costante disequilibrio, mantenuto a spese di energia metabolica fornita dalla molecola di ATP.
\end{itemize}

% ---------------------------------------------------------------------------
\subsection{Definizione e genesi}

\textbf{Potenziale d'azione (PA)}: segnale utilizzato dal sistema nervoso per ricevere, analizzare e trasmettere informazioni. Nasce quando uno stimolo depolarizzante porta il potenziale intracellulare dalla condizione di riposo al \textbf{valore soglia}.

\rubrica{Eccitabilità cellulare}
I neuroni possiedono canali rapidi voltaggio-dipendenti che si aprono in risposta a uno stimolo $\rightarrow$ movimento di ioni attraverso la membrana:
\begin{itemize}
  \item ioni $Na^+$, $Cl^-$ e $Ca^{2+}$ entrano nella cellula;
  \item ioni $K^+$ escono dalla cellula.
\end{itemize}
Il movimento netto di cariche elettriche attraverso la membrana produce una variazione del potenziale di membrana, generando un segnale elettrico.

% ---------------------------------------------------------------------------
\subsection{Fasi del potenziale d'azione}

\rubrica{Fase di riposo}
\begin{itemize}
  \item la permeabilità di membrana agli ioni $K^+$ è alta: i canali sono quasi tutti aperti, in modo che gli ioni $K^+$ possano diffondere fuori dalla cellula;
  \item la permeabilità al $Na^+$ è bassa: il cancello di attivazione è chiuso e quello di inattivazione aperto, e il $Na^+$ non entra nella cellula.
\end{itemize}

\rubrica{Depolarizzazione (ingresso di \texorpdfstring{$Na^+$}{Na+})}
\begin{itemize}
  \item la stimolazione apre il cancello di attivazione dei canali per il $Na^+$;
  \item mentre la cellula si depolarizza, un numero sempre maggiore di canali per il $Na^+$ si apre: la polarità della cellula si inverte (\textbf{overshoot}) fino a $+20$\,mV.
\end{itemize}

\rubrica{Ripolarizzazione (uscita di \texorpdfstring{$K^+$}{K+})}
\begin{itemize}
  \item il cancello di attivazione si chiude e l'ingresso di $Na^+$ cessa;
  \item i più lenti canali del $K^+$ raggiungono il picco di permeabilità.
\end{itemize}

\rubrica{Tipi di variazione del potenziale}
\begin{itemize}
  \item \textbf{depolarizzazione}: il movimento degli ioni rende l'interno della cellula meno negativo o addirittura positivo rispetto all'esterno $\rightarrow$ ingresso di $Na^+$;
  \item \textbf{ripolarizzazione}: il movimento di cariche tende a ristabilire l'originario potenziale di membrana $\rightarrow$ uscita di $K^+$;
  \item \textbf{iperpolarizzazione}: il movimento degli ioni porta il potenziale di membrana a valori più negativi $\rightarrow$ ingresso di $Cl^-$.
\end{itemize}

% ---------------------------------------------------------------------------
\subsection{Propagazione lungo l'assone}

Il potenziale d'azione che origina dalla zona \textit{trigger} del neurone si propaga lungo l'assone; ciò che si propaga è la modificazione della permeabilità della membrana. Il PA è capace di autopropagarsi per l'intera lunghezza della fibra nervosa.

\rubrica{Fibra amielinica}
La conduzione è \textbf{continua}: le fasi di riposo, eccitazione e conduzione dell'impulso si succedono lungo tratti contigui di membrana.

\rubrica{Fibra mielinica}
La conduzione è \textbf{saltatoria}: per la presenza della guaina mielinica il potenziale d'azione nasce solo in corrispondenza dei \textbf{nodi di Ranvier} (zone di membrana in contatto con l'ambiente extracellulare, dove sono possibili gli scambi ionici).

\rubrica{Ruolo della mielina}
\begin{itemize}
  \item riduce la capacità di membrana, ossia la quantità di carica da spostare;
  \item aumenta molto la resistenza di membrana: grazie all'isolamento elettrico della guaina viene persa una quantità minore di segnale;
  \item le fibre mieliniche sono metabolicamente più efficienti, perché il lavoro di pompa è confinato ai nodi di Ranvier, dove sono concentrati i canali del $Na^+$ voltaggio-dipendenti.
\end{itemize}

% ---------------------------------------------------------------------------
\subsection{Le sinapsi}

\textbf{Sinapsi}: struttura deputata a trasmettere il messaggio nervoso tra cellule.
\begin{itemize}
  \item non esiste continuità citoplasmatica tra cellule nervose;
  \item la \textbf{fessura sinaptica} separa la terminazione presinaptica dalla cellula postsinaptica.
\end{itemize}
I neuroni entrano in rapporto tra loro (sinapsi con altri neuroni) o con gli effettori periferici (giunzioni neuromuscolari, giunzioni neuroghiandolari) attraverso contatti sinaptici.

\rubrica{Classificazione morfologica}
\begin{itemize}
  \item \textbf{asso-dendritica};
  \item \textbf{asso-somatica};
  \item \textbf{asso-assonica}.
\end{itemize}

\rubrica{Classificazione funzionale}
\begin{itemize}
  \item \textbf{sinapsi elettriche}: la presenza di \textit{gap junctions} crea una continuità tra le due cellule $\rightarrow$ trasmettono un segnale elettrico direttamente da una cellula all'altra;
  \item \textbf{sinapsi chimiche}: la trasmissione è elettrica-chimica-elettrica $\rightarrow$ rilascio di un neurotrasmettitore.
\end{itemize}

\rubrica{Sinapsi elettrica}
Le \textit{gap junction} poste tra le membrane pre- e post-sinaptiche consentono la comunicazione tra cellule per passaggio diretto passivo di correnti da una cellula all'altra, causando o inibendo la nascita di potenziali d'azione postsinaptici.
\begin{itemize}
  \item si realizza in entrambe le direzioni;
  \item esistono sinapsi elettriche che conducono preferenzialmente in una direzione piuttosto che nell'altra $\rightarrow$ \textbf{rettificazione};
  \item \textbf{trasmissione rapida}: adatta per riflessi in cui sia necessaria una rapida trasmissione tra cellule, ossia una risposta sincrona da parte di un numero elevato di neuroni.
\end{itemize}

\rubrica{Sinapsi chimica}
All'arrivo del potenziale d'azione si attivano i canali per il $Ca^{2+}$; il neurotrasmettitore contenuto nelle vescicole viene rilasciato e attiva sulla cellula post-sinaptica i canali per il $Na^+$, generando un nuovo potenziale d'azione.

% ---------------------------------------------------------------------------
\subsection{\texorpdfstring{Corrente di placca (\textit{End Plate Potential}, EPP)}{Corrente di placca (End Plate Potential, EPP)}}

\rubrica{Versante presinaptico}
La depolarizzazione che si produce all'arrivo del potenziale d'azione determina l'ingresso di $Ca^{2+}$ nel terminale presinaptico, attraverso canali voltaggio-dipendenti che si aprono brevemente, con conseguente esocitosi del mediatore (\textbf{acetilcolina}, sintetizzata da acetil-CoA e colina).

\rubrica{Versante postsinaptico}
L'acetilcolina (ACh) rilasciata nello spazio sinaptico si lega al \textbf{recettore nicotinico}, determinando l'attivazione di un canale cationico che permette l'ingresso di $Na^+$ e $Ca^{2+}$ e l'uscita di $K^+$. La corrente risultante prende il nome di \textbf{corrente di placca} e determina a sua volta una variazione del potenziale in senso depolarizzante (eccitatorio), nota con il nome di \textbf{EPP}.

\rubrica{Rimozione del neurotrasmettitore}
Dopo che il neurotrasmettitore ha svolto la sua azione legandosi al recettore specifico, deve essere rapidamente rimosso per estinguere l'effetto. Per l'ACh questo avviene per degradazione ad opera di enzimi specifici (\textbf{colinesterasi}); la colina viene poi recuperata dal terminale presinaptico.
